\scanpage{29}
Comparons cette équation fonctionnelle avec celle pour la fonction $L_E$, on trouve, en \deleted{\unsure{substituant}} posant pour simplifier $p^{n+\rho}=q$, et $A(t)\delta(E)(-t)^{\chi(E)}=B(t)$ :
\[
L\left(\frac1{qt}\right)=B(t)L(t)
\]
et faisant $t\mapsto\frac1{qt}$ dedans, on trouve
\[
L(t)=B\left(\frac1{qt}\right)L\left(\frac1{qt}\right)=B\left(\frac1{qt}\right)B(t)L(t)
\]
d'où
\[
B(t)B\left(\frac1{qt}\right)=1.
\]
Remplaçant $B$ par sa valeur, on trouve
\[
L_E^\infty\left(\frac1{\deleted{\unsure{p^{n+\rho}}}qt}\right)^{-1}\delta(E)(-t)^{\chi(E)}L_E^\infty(t)^{-1}\delta(E)\left(\frac{-1}{qt}\right)^{\chi(E)}=1
\]
i.e.
\[
\deleted{\unsure{L_E^\infty(\frac1{p^{n+\rho}t})^{-1}|L_E^\infty(p^{n+\rho}t)}}
\qquad L_E^\infty(t)L_E^\infty\left(\frac1{qt}\right)=\delta(E)^2\deleted{\frac1q}q^{-\chi(E)}.
\]
C'est précisément l'équation fonctionnelle explicitée plus haut.

Posons
\[
\unsure{\xi|\zeta}_E(t)=t^{\chi(E)/2}L_E(t)
\]
alors
\[
\begin{aligned}
\unsure{\xi|\zeta}_E\left(\frac1{qt}\right)&=\left(\frac1{qt}\right)^{\chi(E)/2}L_E\left(\frac1{qt}\right)=\left(\frac1{qt}\right)^{\chi(E)/2}A(t)\delta(E)(-t)^{\chi(E)}L_E(t)\\
&=\underbrace{q^{-\chi(E)/2}A(t)\delta(E)}(-1)^{\chi(E)}\underbrace{t^{\chi(E)/2}L_E(t)}_{\unsure{\xi|\zeta}_E(t)}
\end{aligned}
\]
i.e.
\mspage{28}
